\section{Limitations}

\textbf{Exposure is not causation.} Our blast-radius metrics measure
whether a memory is reachable from a compromised root under a given
provenance policy and whether it actually carries the harmful semantic
content---they do not establish that the exposure caused any real-world
harm (an action taken, a decision made). A memory can be correctly
flagged as exposed and contaminated without ever being acted upon.

\textbf{Synthetic, LLM-authored corpus.} All 30 scenarios are
hand-authored specifications with LLM-generated surface text, not
derived from real production agent-memory traces or naturally
occurring attacks. We deliberately designed a real-document validation
slice into the original study plan (10--20 real documents, for
ecological validity) specifically to partially address this; it was
not completed within this paper's timeline and is left as a genuine,
acknowledged gap rather than a claim we did not intend to make.

\textbf{Single embedding model, never pinned as the paper's actual
choice.} Retrieval and the H2 attribution-score analysis both use
\texttt{all-MiniLM-L6-v2} (sentence-transformers) throughout---explicitly
flagged in our own design documents as a placeholder pending a final
pinned choice, not yet revisited before this draft. Results involving
embedding similarity (H2's threshold sweep) should be read with this
caveat; a different embedding model could shift the precise
precision/recall frontier, though we would not expect it to change the
qualitative existence of a frontier.

\textbf{\texttt{attribution\_threshold} was operationalized after the
fact, not specified in the original design.} The experimental grid
named this analysis factor without defining a continuous per-edge score
to threshold. We chose cosine similarity between parent/child memory
embeddings---an explicit, defensible, but not the only possible choice
(a trained classifier or NLI-entailment-based score would also have
been reasonable)---documented as a dated decision, not silently
assumed. The same applies to the \texttt{depth\_aware} containment
policy, which the design named as ``one baseline, one proposed method''
without specifying the proposed method's algorithm; we operationalized
it as a depth-windowed hybrid of conservative and structural-only
propagation. Both choices are reasoned and stated, but a different
operationalization of either could produce a different quantitative
frontier. We accordingly present \texttt{depth\_aware} as a
proof-of-concept operating point demonstrating that depth-dependent
provenance decisions occupy a useful position on the quarantine/recall
frontier --- not as a proposed, tuned, or universally optimal
containment algorithm.

\textbf{Conservative propagation gives recall 1.0 largely by
construction.} The context-exposure/\texttt{flat\_transitive} policy's
near-perfect recall ($R_{BR} \approx 1.0$ across nearly all conditions)
is expected given how it is defined (flag anything reachable via any
edge), not a discovery---the paper's contribution is characterizing the
precision \emph{cost} of that recall and the conditions under which
cheaper (structural-only, thresholded, depth-aware) policies still miss
real contamination, not the recall-1.0 fact itself.

\textbf{One containment-cost metric turned out uninformative in our
topology.} Our original operationalization of $C_{regen}$ (distinct
generating runs touched) does not differentiate \texttt{flat\_transitive}
from \texttt{depth\_aware} in this harness's design, because the
always-continuing structural lineage guarantees every depth is touched
under either policy. We report this honestly (Method, Results) and use
flagged \emph{object count} as the metric that actually demonstrates
the containment tradeoff instead---but it means our depth-based cost
metric, as originally specified, does not do the job we designed it to
do, and a different harness topology (e.g.\ one where structural
lineage can itself terminate early) might make it more informative.

\textbf{Depth-continuation rule limits branching complexity.} Only the
compromised root's own lineage (\texttt{child\_1}) continues across
depths; sibling writes are fresh at each depth rather than themselves
branching further. This keeps the corpus tractable and auditable but
means we do not study \emph{nested} branching (a sibling write itself
spawning further descendants)---a real agent's memory graph could have
branching at every node, not just along one continuing spine.

\textbf{Three models, one derivation prompt template.} We test three
LLMs (gpt-4o-mini, gpt-4o, Llama-3.3-70B-Instruct) with an identical,
un-tuned prompt template, by design---this isolates the model as the
variable of interest, but it also means results may not reflect each
model's \emph{best-achievable} derivation behavior under a prompt
specifically tuned for it, and does not cover reasoning-specialized or
much smaller/larger models.

\textbf{Platform-level content filtering caused incomplete coverage
for one model.} 30 of 21,600 planned traces (Llama-3.3-70B-Instruct
only, 0.14\% of that model's grid) could not be generated because
Azure's content-safety filter rejected the underlying request (flagged
as ``Jailbreak'')---a platform behavior, not a code or methodology
defect, but one that means Llama-3.3-70B's corpus has a small number of
missing cells the other two models' corpora do not, concentrated in
\texttt{authoritative\_framing}-style scenarios and
\texttt{write\_fanout=3} conditions. Left as genuine missing data, not
imputed.

\textbf{Human validation now includes a human-human baseline.} Cohen's
$\kappa$ between the automated labeler and one human's blind labels on
a 120-sample calibration set is 0.87--0.88. A second independent human
annotator blind-labeled the same 120-sample set afterward, giving a
human-human baseline of $\kappa = 0.9277$ (115/120 raw agreement) ---
close to the pipeline's own agreement with either human, and
informative rather than merely reassuring: all 5 human-human
disagreements fall on the same REFERENCES-vs-CARRIES compositional
boundary already flagged below as the labeler's weakest point,
indicating that boundary reflects genuine rubric ambiguity rather than
a labeler-specific failure (see Discussion).

\textbf{Marker-token extraction for H3 is mechanical, not
human-authored.} No scenario specification includes an explicit
\texttt{marker\_tokens} field; we extract them via a regex-plus-entity-substring
heuristic verified to produce at least one marker per scenario, but not
independently reviewed the way the scenario corpus's structural ground
truth was.
